\documentclass[12pt,a4paper]{article}

\usepackage[a4paper,margin=1in]{geometry}
\usepackage{amsmath,amssymb,mathtools}
\usepackage{xcolor}
\usepackage{hyperref}
\usepackage{array}
\usepackage{booktabs}
\usepackage{enumitem}

\hypersetup{
    colorlinks=true,
    linkcolor=blue,
    urlcolor=blue
}

\setlength{\parindent}{0pt}
\setlength{\parskip}{6pt}

% 自定义微分符号宏定义
\newcommand{\dif}{\mathop{}\!\mathrm{d}}

\begin{document}

\begin{center}
    {\color[HTML]{2E86C1}\fontsize{24}{28}\selectfont\textbf{LaTeX Workshop Worksheet 2 -- Formula}}
\end{center}

Welcome to the {\color[HTML]{117A65}\textbf{LaTeX Workshop Formula Part}}.

{\color[HTML]{7D3C98}\textbf{Goal:}} Learn the basics of LaTeX math, improve clarity and consistency, and practice writing formulas using professional conventions.

\hrulefill

\section{\color[HTML]{C0392B}Common Formulas (Quick Reference)}

Here are some \textbf{commonly used formulas} in mathematics, already written in LaTeX.
Use them as examples and modify them in the exercises.

Remember to add \verb|\usepackage{amsmath}| at the beginning of your code please!

\subsection{Fractions \& Roots}

\verb|\[\frac{a+b}{c+d}, \quad \sqrt{x^2+y^2}, \quad \sqrt[3]{8}=2\]|

\textbf{After compilation}

\[
\frac{a+b}{c+d}, \quad \sqrt{x^2+y^2}, \quad \sqrt[3]{8}=2
\]

\textbf{\color[HTML]{CA6F1E}Exercises}

Write the quadratic formula:
\[
x=\frac{-b\pm\sqrt{b^2-4ac}}{2a}
\]

\subsection{Exponents \& Logarithms}

\begin{verbatim}
\[
e^2 \approx 7.389, \quad e^{i\pi}+1=0, \quad
\log_a b = \frac{\ln b}{\ln a}
\]
\end{verbatim}

\textbf{After compilation}

\[
e^2 \approx 7.389, \quad e^{i\pi}+1=0, \quad
\log_a b = \frac{\ln b}{\ln a}
\]

\subsection{Limits}

\begin{verbatim}
\[
\lim_{x\to0}\frac{\sin x}{x}=1, \quad
\lim_{n\to\infty}\left(1+\frac{1}{n}\right)^n=e
\]
\end{verbatim}

\textbf{After compilation}

\[
\lim_{x\to0}\frac{\sin x}{x}=1, \quad
\lim_{n\to\infty}\left(1+\frac{1}{n}\right)^n=e
\]

\textbf{\color[HTML]{CA6F1E}Exercises}

Write the classic limit:
\[
\lim_{x\to0}\frac{\sin x}{x}=1
\]

\subsection{Derivatives}

\begin{verbatim}
\[
\frac{d}{dx}\sin x=\cos x, \quad
f'(x)=\lim_{h\to0}\frac{f(x+h)-f(x)}{h}
\]    
\end{verbatim}

\textbf{After compilation}

\[
\frac{d}{dx}\sin x=\cos x, \quad
f'(x)=\lim_{h\to0}\frac{f(x+h)-f(x)}{h}
\]

\subsection{Integrals}

\begin{verbatim}
\[
\int_0^1 x^2\,dx=\frac{1}{3}, \quad
\int_{-\infty}^{\infty}e^{-x^2}\,dx=\sqrt{\pi}
\]
\end{verbatim}

\textbf{After compilation}

\[
\int_0^1 x^2\,dx=\frac{1}{3}, \quad
\int_{-\infty}^{\infty}e^{-x^2}\,dx=\sqrt{\pi}
\]

\textbf{\color[HTML]{CA6F1E}Exercises}

Write the integral:
\[
\int_0^\infty e^{-x}\,dx=1
\]

\subsection{Summations \& Products}

\begin{verbatim}
\[
\sum_{i=1}^{n}i=\frac{n(n+1)}{2}, \quad
\prod_{k=1}^{n}k=n!
\]
\end{verbatim}

\textbf{After compilation}

\[
\sum_{i=1}^{n}i=\frac{n(n+1)}{2}, \quad
\prod_{k=1}^{n}k=n!
\]

\textbf{\color[HTML]{CA6F1E}Exercises}

Write the summation:
\[
\sum_{n=1}^{\infty}\frac{1}{n^2}=\frac{\pi^2}{6}
\]

\subsection{Vectors \& Matrices}

\begin{verbatim}
\[
\vec{v}=\begin{bmatrix}x\\y\\z\end{bmatrix},\quad
A=\begin{Bmatrix}1&0\\0&1\end{Bmatrix},\quad
B=\begin{pmatrix}1&2\\3&4\end{pmatrix},\quad
C=\begin{vmatrix}1&2\\3&4\end{vmatrix}
\]
\end{verbatim}


\textbf{After compilation}

\[
\vec{v}=\begin{bmatrix}x\\y\\z\end{bmatrix},\quad
A=\begin{Bmatrix}1&0\\0&1\end{Bmatrix},\quad
B=\begin{pmatrix}1&2\\3&4\end{pmatrix},\quad
C=\begin{vmatrix}1&2\\3&4\end{vmatrix}
\]

\textbf{\color[HTML]{CA6F1E}Exercises}

Write the vector and the matrix:
\[
\vec{v}=(x,y,z),\quad 
A=
\begin{bmatrix}
1&0&0\\
0&1&0\\
0&0&1
\end{bmatrix}
\]

\subsection{Set Operator}

\begin{verbatim}
\[
< \quad > \quad \subset \quad \supset \quad
\subseteq \quad \supseteq \quad \cup \quad \cap
\]
\end{verbatim}

\textbf{After compilation}

\[
< \quad > \quad \subset \quad \supset \quad
\subseteq \quad \supseteq \quad \cup \quad \cap
\]

\subsection{Special Symbols}

\begin{verbatim}
\[
\alpha,\ \beta,\ \gamma,\ \delta,\ \theta,\ \infty,\ 
\forall,\ \exists,\ \nabla
\]
\end{verbatim}

\textbf{After compilation}

\[
\alpha,\ \beta,\ \gamma,\ \delta,\ \theta,\ \infty,\ 
\forall,\ \exists,\ \nabla
\]

\textbf{Tips}

\begin{enumerate}
    \item Check \href{https://www.overleaf.com/learn/latex/Mathematical_expressions}{Math expressions} for more detailed information.
    \item Use \href{https://www.latexlive.com/}{Formula editor} to generate formula automatically.
    \item Smart use of AI!
\end{enumerate}

\section{\color[HTML]{C0392B}Math Modes in LaTeX}

LaTeX provides \textbf{different ways} to enter math mode.
It is important to use the right one depending on whether you want inline or display math.

\subsection{Inline Math}

\begin{itemize}
    \item \textbf{Syntax:} \verb|\(...\)| or single dollar \verb|$...$|
    \item \textbf{Use:} Inside a sentence.
    \item \textbf{Example:}
\end{itemize}

\begin{verbatim}
This is inline math: \( a^2 + b^2 = c^2 \).
Or equivalently: $a^2+b^2=c^2$.
\end{verbatim}

\textbf{After compilation}

This is inline math: \( a^2 + b^2 = c^2 \).
Or equivalently: $a^2+b^2=c^2$.

In professional LaTeX (academic journals), \verb|\(...\)| is preferred because \verb|$...$| is TeX legacy and may cause spacing issues.

\subsection{Display Math}

\begin{itemize}
    \item \textbf{Syntax:} \verb|\[...\]| or double dollars \verb|$$...$$|
    \item \textbf{Use:} Standalone equation, centered, without numbering.
    \item \textbf{Example:}
\end{itemize}

\begin{verbatim}
The mass-energy equation below reveals the equivalence relationship
between mass and energy: \[ E = mc^2 \]
\end{verbatim}

\textbf{After compilation}

The mass-energy equation below reveals the equivalence relationship
between mass and energy: \[ E = mc^2 \]

\subsection{Numbered Equations}

\begin{itemize}
    \item \textbf{Syntax:} \texttt{equation} environment
    \item \textbf{Use:} To number a formula automatically.
    \item \textbf{Example:}
\end{itemize}

\begin{verbatim}
\begin{equation}
\nabla \cdot \vec{E} = \frac{\rho}{\varepsilon_0}
\end{equation}

\begin{equation}
E = mc^2
\end{equation}
\end{verbatim}

\textbf{After compilation}

\begin{equation}
\nabla \cdot \vec{E} = \frac{\rho}{\varepsilon_0}
\end{equation}

\begin{equation}
E = mc^2
\end{equation}

\subsection{Multi-line Alignment}

\begin{itemize}
    \item \textbf{Syntax:} \texttt{align} (from \textbf{amsmath})
    \item \textbf{Use:} Align multi-step derivations at $=$ or other operators.
    \item \textbf{Example:}
\end{itemize}

\begin{verbatim}
\begin{align}
f(x) &= (x+1)^2 \\
     &= x^2 + 2x + 1
\end{align}

\begin{align}
f(x) &= (x+1)^2 \notag \\
     &= x^2 + 2x + 1
\end{align}

\begin{equation}
\begin{aligned}
f(x) &= (x+1)^2 \\
     &= x^2 + 2x + 1
\end{aligned}
\end{equation}
\end{verbatim}

\textbf{After compilation}
\begin{align}
f(x) &= (x+1)^2 \\
     &= x^2 + 2x + 1
\end{align}

\begin{align}
f(x) &= (x+1)^2 \notag \\
     &= x^2 + 2x + 1
\end{align}

\begin{equation}
%\notag
\begin{aligned}
f(x) &= (x+1)^2 \\
     &= x^2 + 2x + 1
\end{aligned}
\end{equation}

\textbf{\color[HTML]{CA6F1E}Exercises}

Use the \textbf{align environment} to typeset

\begin{align}
f(x)&=(x+1)^3\\
&=x^3+3x^2+3x+1
\end{align}

\section{\color[HTML]{C0392B}Best Reminder \& Tips}

\subsection{Multi-line Alignment: \texttt{align} vs \texttt{aligned}}

\paragraph{align}

\begin{itemize}
    \item Requires \texttt{amsmath} package: \verb|\usepackage{amsmath}|
    \item Standalone multi-line equations with automatic numbering
    \item Example:
\end{itemize}

\begin{verbatim}
\begin{align}
f(x) &= (x+1)^2 \\
     &= x^2 + 2x + 1
\end{align}
\end{verbatim}

\textbf{Notes:}
\begin{itemize}
    \item Use \verb|&| to indicate alignment points (typically at \verb|=|)
    \item Use \verb|\nonumber| or \verb|\notag| to remove numbering on specific lines
    \item Use \verb|\label{}| for cross-references with \verb|\eqref{}|
\end{itemize}

\paragraph{aligned}
\begin{itemize}
    \item Also requires \texttt{amsmath}, but must be inside another math environment
    \item Does not provide numbering on its own
    \item Example:
\end{itemize}

\begin{verbatim}
\[
\begin{aligned}
f(x) &= (x+1)^2 \\
     &= x^2 + 2x + 1
\end{aligned}
\]
\end{verbatim}

\textbf{Notes:}
\begin{itemize}
    \item Only allowed in math environment. (Need \verb|\begin{equation}| or \verb|\[...\]|)
    \item Ideal for embedding aligned equations inline or inside display math
    \item For numbering, wrap in \texttt{equation} environment:
\end{itemize}

\begin{verbatim}
\begin{equation}
\notag
\begin{aligned}
f(x) &= (x+1)^2 \\
     &= x^2 + 2x + 1
\end{aligned}
\end{equation}
\end{verbatim}

\subsection{Use \texttt{\textbackslash mathrm\{d\}x} for Differentials}

\begin{itemize}
    \item Differential $d$ is an operator, not a variable; it should be upright.
    \item Example:
\end{itemize}

\[
\int_0^1 x^2\,\mathrm{d}x
\]

\begin{verbatim}
\[\int_0^1 x^2 \, \mathrm{d}x\]
\end{verbatim}

\textbf{Optional macro:}

\[
\int_0^1 x^2\,\dif x
\]

\begin{verbatim}
\newcommand{\dif}{\mathop{}\!\mathrm{d}}
\[\int_0^1 x^2\,\dif x\]
\end{verbatim}

\subsection{Control Spacing}

\begin{center}
\begin{tabular}{c l l}
\toprule
Command & Width & Example \\
\midrule
\verb|\,| & thin space (1/6 quad) & $a\,b$ \\
\verb|\:| & medium space (2/9 quad) & $a\:b$ \\
\verb|\;| & thick space (5/18 quad) & $a\;b$ \\
\verb|\!| & negative thin space (-1/6 quad) & $a\!b$ \\
\verb|\quad| & 1 em (width of ``M'') & $a\quad b$ \\
\verb|\qquad| & 2 em & $a\qquad b$ \\
\bottomrule
\end{tabular}
\end{center}

\textbf{1. Multiplication}

Prefer \verb|\cdot| or thin space \verb|\,| over implicit juxtaposition when variables are next to numbers.

\[
2\,x \quad \text{or} \quad 2\cdot x
\]

\begin{verbatim}
\[
2\,x \quad \text{or} \quad 2\cdot x
\]
\end{verbatim}

\textbf{2. Function Application}

Always put a thin space after functions before arguments for clarity.

\[
\sin\,x,\quad \log\,n
\]

\begin{verbatim}
\[
\sin\,x,\quad \log\,n
\]
\end{verbatim}

\textbf{3. Fractions}

Use small spaces around \verb|\frac| if needed for readability.

\[
\frac{a}{b}\,x
\]

\begin{verbatim}
\[
\frac{a}{b}\,x
\]
\end{verbatim}

\textbf{4. Operators in Limits/Integrals}

Add thin or medium space around $d$ in integrals.

\[
\int_0^1 f(x)\,\mathrm{d}x
\]

\begin{verbatim}
\[
\int_0^1 f(x)\,\mathrm{d}x
\]
\end{verbatim}

\textbf{5. Negative Space}

Use \verb|\!| sparingly to reduce awkward gaps in products or superscripts.

\[
x^{n\!+1}
\]

\begin{verbatim}
\[
x^{n\!+1}
\]
\end{verbatim}

To see more detailed information in section 7 of this webpage:
\href{https://latex-tutorial.com/latex-space/}{latex\_space}.

\subsection{Place Limits Explicitly}

\begin{itemize}
    \item Inline math usually puts limits to the right:\, \(\sum_{i=1}^{n}i\)
    
    \item Use \verb|\limits| to force display-style limits inline:
    \(\sum\limits_{i=1}^{n}i\)
    
    \item Or \verb|\displaystyle|:
    \(\displaystyle\sum_{i=1}^{n}i\)
\end{itemize}

\begin{verbatim}
\(\sum\limits_{i=1}^{n}i\)
\end{verbatim}

\begin{verbatim}
\(\displaystyle\sum_{i=1}^{n}i\)
\end{verbatim}

\subsection{Use \texttt{\textbackslash label} and \texttt{\textbackslash ref} for Numbering}

\begin{itemize}
    \item Automatic numbering and cross-references
    \item Example:
\end{itemize}

\begin{equation}
\label{eq:problem1_quadratic}
x = \frac{-b \pm \sqrt{b^2 - 4ac}}{2a}
\end{equation}

As shown in equation \eqref{eq:problem1_quadratic}, ...

\begin{verbatim}
\begin{equation}
\label{eq:problem1_quadratic}
x = \frac{-b \pm \sqrt{b^2 - 4ac}}{2a}
\end{equation}

As shown in equation \eqref{eq:problem1_quadratic}, ...
\end{verbatim}

\begin{itemize}
    \item \verb|\ref{eq:problem1_quadratic}| $\rightarrow$ prints just the number: 1.
    \item \verb|\eqref{eq:problem1_quadratic}| $\rightarrow$ prints the number with parentheses: (1).
\end{itemize}

\section{\color[HTML]{2874A6}Other Related Materials}

\begin{enumerate}
    \item \href{https://www.tablesgenerator.com/}{AutoTableGenerate}
    \item \href{https://www.ctan.org/}{The Comprehensive TeX Archive Network}
\end{enumerate}

\end{document}